\chapter{Grafy rozkładów JSJ i grupy trójrozmaitości}
\label{ch:tg-jsj-graphs}

Graf, którego wierzchołki opisują części rozmaitości,
a krawędzie powierzchnie ich sklejania, jest jednym z
najważniejszych modeli współczesnej topologii
trójwymiarowej. Po sekcji~\ref{sec:kt-dehn}, poświęconej
chirurgii Dehna, możemy go zastosować do zewnętrzy węzłów.
Odróżniamy \emph{graf dualny rozkładu}
od całej \emph{struktury grafu przestrzeni lub grup}:
sam graf zapomina, jakie rozmaitości sklejamy i jak
identyfikujemy ich brzegi. Różnica jest analogiczna
do tej, którą poznaliśmy przy triangulacjach powierzchni.

\section{Torusy nieściśliwe i definicja grafu JSJ}
\label{sec:tg-jsj-graph-def}

\begin{definicja}[Nieściśliwy torus]
Niech $M$ będzie zwartą, spójną, orientowalną
trójrozmaitością, a $T\subset\operatorname{int}M$
dwustronnym, gładko zanurzonym torusem.
Przyjmujemy kategorię gładką, aby wykluczyć dzikie
zanurzenia topologiczne. Zgodnie z sekcją
\ref{sec:kt-torus-jsj} nazywamy $T$ \emph{nieściśliwym},
gdy indukowany inkluzją homomorfizm
$\pi_1(T)\to\pi_1(M)$ jest injekcją.
Torus $T$ nazywamy \emph{równoległym do brzegu},
gdy ogranicza wraz z pewną składową $\partial M$
podrozmaitość homeomorficzną z $T^2\times[0,1]$.
\end{definicja}

\begin{twierdzenie}[Jaco--Shalen--Johannson; rozkład JSJ]
\label{thm:tg-jsj-theorem}
Niech $M$ będzie zwartą, spójną, orientowalną,
nieprzywiedlną trójrozmaitością, której brzeg jest
pusty lub jest sumą torusów. Istnieje skończona
rodzina parami rozłącznych, nieściśliwych,
nierównoległych do brzegu torusów, parami
nieizotopowych, taka że każda składowa rozmaitości
rozciętej wzdłuż tych torusów jest rozmaitością
Seiferta lub atoroidalną. Wśród rodzin o tej
własności rodzina o najmniejszej liczbie torusów jest wyznaczona
jednoznacznie z dokładnością do izotopii całej rodziny.
Dopuszczamy rodzinę pustą.
\end{twierdzenie}

\begin{uwaga}[Status twierdzenia]
Nieprzywiedlność oznacza, że każda gładko zanurzona
sfera $S^2\subset\operatorname{int}M$ ogranicza w $M$
zanurzoną kulę trójwymiarową. Rozmaitość geometrycznie atoroidalna
to taka, w której każdy nieściśliwy torus jest
równoległy do brzegu. Atoroidalność homotopijna
oznacza natomiast, że każde odwzorowanie $T^2\to M$
indukujące monomorfizm na grupach podstawowych
można homotopować do odwzorowania o obrazie w brzegu.
Pojęcia te nie zawsze są równoważne (wyjątki
wśród małych rozmaitości Seiferta). W twierdzeniu
JSJ używamy warunku geometrycznego; por. dyskusję
przed twierdzeniem 1.6.1 monografii
Aschenbrennera--Friedla--Wiltona. Rozmaitość Seiferta posiada
rozwarstwienie na okręgi z dopuszczalnymi włóknami
wyjątkowymi. Samo istnienie i jednoznaczność
minimalnego rozkładu JSJ przyjmujemy jako głęboki
wynik topologii trójwymiarowej; zob.
M. Aschenbrenner, S. Friedl, H. Wilton,
\emph{3-Manifold Groups}, EMS, 2015, twierdzenie 1.6.1
i poprzedzającą je dyskusję o definicjach atoroidalności.
Dla rozmaitości już atoroidalnej lub będącej
rozmaitością Seiferta rodzina JSJ jest pusta;
istnienie torusa nieściśliwego nie wymusza więc
niepustego rozkładu JSJ. Nie należy twierdzenia
wywodzić jedynie z klasycznego twierdzenia
van Kampena.
\end{uwaga}

\begin{definicja}[Graf dualny rozkładu JSJ]
Niech $\mathcal T=\{T_1,\ldots,T_m\}$ będzie
rodziną torusów o najmniejszej liczności z twierdzenia
\ref{thm:tg-jsj-theorem},
a $M\mid\mathcal T$ rozmaitością z brzegiem,
otrzymaną przez rozcięcie $M$ wzdłuż tych torusów.
\emph{Graf JSJ} $\Gamma_{\mathrm{JSJ}}(M)$ ma po jednym
wierzchołku dla każdej składowej $M\mid\mathcal T$
oraz po jednej krawędzi dla każdego $T_i$;
końce krawędzi odpowiadają składowym przylegającym
do dwóch stron $T_i$. Jeżeli obie strony leżą
w jednej składowej, otrzymujemy pętlę. Dwa różne
torusy mogą wyznaczać krawędzie o tych samych końcach,
więc graf rozumiemy jako multigraf z dopuszczalnymi
pętlami. Brzegi pierwotnej rozmaitości $M$ nie
wyznaczają krawędzi: nie są powierzchniami sklejania.
\end{definicja}

\begin{stwierdzenie}[Spójność grafu rozkładu]
\label{st:tg-jsj-connected}
Dla każdej rozmaitości $M$ spełniającej założenia
poprzedniej definicji graf $\Gamma_{\mathrm{JSJ}}(M)$
jest spójny. Gdy $m=0$, ma dokładnie jeden
wierzchołek i żadnej krawędzi.
\end{stwierdzenie}
\begin{proof}
Niech $M_i,M_j$ będą dwiema składowymi rozciętej
rozmaitości. Wybieramy punkty $p_i\in\operatorname{int}M_i$
i $p_j\in\operatorname{int}M_j$. Ponieważ $M$ jest
spójna i łukowo spójna, istnieje ścieżka $\gamma$
łącząca ich obrazy w $M$. Wybieramy gładką ścieżkę we wnętrzu $M$ (kołnierz
brzegu pozwala odsunąć ścieżkę od brzegu).
Przez twierdzenie o transwersalności możemy ją
perturbować z ustalonymi końcami tak, aby była
poprzeczna do każdego torusa. Przeciwobraz każdego
torusa jest wtedy zwartą podrozmaitością wymiaru
zero przedziału, a więc jest skończony. Kolejne fragmenty
ścieżki leżą w składowych rozciętej rozmaitości,
a każde przecięcie torusa odpowiada krawędzi grafu.
Otrzymujemy spacer z $M_i$ do $M_j$, co dowodzi
spójności. Jeżeli torusów nie ma, rozcięta
rozmaitość pozostaje równa spójnej $M$.
\end{proof}

\begin{stwierdzenie}[Torus rozdzielający i most]
\label{st:tg-jsj-bridge}
Torus $T_e$ z rodziny JSJ rozdziela $M$ wtedy
i tylko wtedy, gdy odpowiadająca mu krawędź $e$
jest mostem w grafie $\Gamma_{\mathrm{JSJ}}(M)$.
\end{stwierdzenie}
\begin{proof}
Rozetnijmy $M$ wzdłuż całej rodziny $\mathcal T$,
a następnie sklejmy z powrotem wszystkie pary
kopii torusów poza parą odpowiadającą $e$.
Otrzymujemy rozmaitość rozciętą tylko wzdłuż $T_e$.
Jej składowe odpowiadają dokładnie składowym grafu
$\Gamma_{\mathrm{JSJ}}(M)-e$: sklejanie wzdłuż
niepustej powierzchni łączy dwie spójne części,
jeżeli dotąd leżały w różnych składowych.
Indukując po liczbie wykonanych sklejeń, uzyskujemy
tę odpowiedniość również dla pętli i krawędzi
równoległych. Torus rozdziela $M$, gdy rozcięcie
jest niespójne. Ponieważ pierwotny graf jest spójny
na mocy stwierdzenia \ref{st:tg-jsj-connected},
jest to dokładnie warunek, że $e$ jest mostem.
\end{proof}

\begin{center}
\begin{tikzpicture}[scale=.95,line cap=round]
 \node[tgBlueVertex] (A) at (-2,0) {};
 \node[tgTealVertex] (B) at (0,.85) {};
 \node[tgRedVertex] (C) at (2,0) {};
 \draw[tgEdge] (A)--node[above left,font=\scriptsize] {$T_1$}(B);
 \draw[tgEdge] (B)--node[above right,font=\scriptsize] {$T_2$}(C);
 \node[font=\scriptsize,text=tgNavy] at (-2,-.62) {$M_1$};
 \node[font=\scriptsize,text=tgNavy] at (0,1.43) {$M_2$};
 \node[font=\scriptsize,text=tgNavy] at (2,-.62) {$M_3$};
 \node[font=\footnotesize,text=tgNavy] at (0,-1.12)
   {Wierzchołki oznaczają części; krawędzie -- torusy sklejania.};
\end{tikzpicture}
\end{center}

\section{Grupa fundamentalna z grafu JSJ}
\label{sec:tg-jsj-groups}

\begin{stwierdzenie}[Rozdzielający torus i amalgamacja]
Niech $M=M_1\cup_T M_2$, gdzie $M_1,M_2$ są
spójnymi trójrozmaitościami, a $T$ jest wspólnym,
spójnym torusem brzegowym. Załóżmy, że inkluzje
$T\hookrightarrow M_i$ indukują monomorfizmy na
$\pi_1$. Wtedy, po zgodnym wyborze punktów i ścieżek
bazowych,
\[
\pi_1(M)\cong\pi_1(M_1)
 *_{\pi_1(T)}\pi_1(M_2),
\qquad \pi_1(T)\cong\mathbb Z^2.
\]
\end{stwierdzenie}
\begin{proof}
Wybieramy kołnierz $T\times(-\varepsilon,\varepsilon)$
i nieco powiększamy $M_1,M_2$ o odpowiednie części
kołnierza, aby powstało pokrycie $M$ dwoma
otwartymi, łukowo spójnymi zbiorami. Ich część
wspólna deformacyjnie retraktuje się na $T$, a
każdy zbiór na odpowiednie $M_i$. W twierdzeniu
van Kampena dostajemy prezentację z generatorami
z obu grup $\pi_1(M_i)$ i relacjami utożsamiającymi
obrazy elementów $\pi_1(T)$. Ponieważ założyliśmy
injektywność obu odwzorowań, otrzymany iloraz
iloczynu wolnego jest dokładnie iloczynem wolnym
z amalgamacją nad $\mathbb Z^2$.
\end{proof}

\begin{uwaga}[Torus nierozdzielający]
Jeśli po rozcięciu wzdłuż nieściśliwego torusa
$T$ otrzymujemy spójną rozmaitość $N$ o dwóch
wyróżnionych składowych toroidalnego brzegu,
to sklejenie ich z powrotem jest dokładnie
przypadkiem HNN z twierdzenia \ref{thm:tg-vk-amalgam-hnn}.
Oba włączenia kopii $T$ do $N$ są injektywne
na grupach fundamentalnych, ponieważ ich złożenia
z naturalnym odwzorowaniem $N\to M$ są (z dokładnością
do sprzężenia) monomorfizmem $\pi_1(T)\to\pi_1(M)$.
Po wyborze dróg bazowych otrzymujemy prezentację
\[
\pi_1(M)\cong
\langle \pi_1(N),t\mid
t\iota_-(a)t^{-1}=\iota_+(a),\ a\in\pi_1(T)\rangle.
\]
Wybieramy orientację stabilnej litery $t$ zgodnie
z konwencją przyjętą w przywołanym twierdzeniu.
Odpowiada ona pętli przechodzącej przez sklejany torus.
Cały graf JSJ, wyposażony w grupy
$\pi_1(M_v)$ i $\pi_1(T_e)$ oraz odpowiednie
monomorfizmy, jest grafem grup. Nie wystarcza
sam nieoznakowany graf $\Gamma_{\mathrm{JSJ}}(M)$.
\end{uwaga}

\begin{twierdzenie}[Grupa fundamentalna rozkładu JSJ]
\label{thm:tg-jsj-graph-of-groups}
Niech $M$ będzie zwartą, spójną, orientowalną,
nieprzywiedlną trójrozmaitością o brzegu pustym
lub będącym sumą torusów. Niech $\mathcal T$
będzie jej rodziną torusów JSJ,
$\Gamma=\Gamma_{\mathrm{JSJ}}(M)$ grafem
dualnym, a $M_v$ składową rozciętej
rozmaitości odpowiadającą wierzchołkowi $v$.
Przypiszmy wierzchołkom grupy
$G_v=\pi_1(M_v)$, a krawędziom
$G_e=\pi_1(T_e)\cong\mathbb Z^2$,
z monomorfizmami wyznaczonymi przez
włączenia torusów do przyległych części
(po wyborze dróg bazowych).
Wtedy
\[
\pi_1(M)\cong\pi_1(\mathcal G),
\]
gdzie $\mathcal G$ jest tak otrzymanym grafem grup.

Jeżeli $\Gamma$ ma $s$ wierzchołków i $m$
krawędzi, to można skonstruować
$\pi_1(M)$, wykonując najpierw
$s-1$ amalgamacji wzdłuż krawędzi
wybranego drzewa rozpinającego $\Gamma$,
a następnie $m-s+1$ rozszerzeń HNN
wzdłuż pozostałych krawędzi.
\end{twierdzenie}

\begin{proof}
Wybierzmy w $\Gamma$ drzewo rozpinające
$T_\Gamma$, którego istnienie wynika
ze stwierdzenia \ref{st:tg-jsj-connected}.
Rozcięcie daje dwie kopie każdego torusa.
Między odpowiadające im brzegi wstawiamy cylinder
$T_e\times[0,1]$, którego końce przyklejamy
zgodnie z pierwotnymi identyfikacjami.
Wstawienie cylindra nie zmienia typu
homeomorficznego, gdyż dwustronny kołnierz
torusa można przeparametryzować. W ten sposób
otrzymujemy graf przestrzeni homeomorficzny z $M$,
z przestrzeniami wierzchołkowymi $M_v$
i cylindrami krawędziowymi.
Ponieważ torusy JSJ są nieściśliwe,
odwzorowania $\pi_1(T_e)\to\pi_1(M_v)$
są injektywne. Jest to konsekwencja
injektywności $\pi_1(T_e)\to\pi_1(M)$
oraz faktoryzacji tego odwzorowania
przez $\pi_1(M_v)$. Ten sam argument obowiązuje
na każdym pośrednim etapie sklejania: dla każdej
jeszcze niesklejonej kopii torusa złożenie jej
włączenia do bieżącej części z naturalnym
odwzorowaniem tej części do $M$ indukuje
monomorfizm na grupach fundamentalnych. Zatem
włączenie do części pośredniej jest również
monomorfizmem. Nie zakładamy tu z góry
injektywności grup wszystkich części w $\pi_1(M)$.

Łączmy części w kolejności
krawędzi drzewa $T_\Gamma$.
Każda kolejna krawędź łączy
dwie dotychczas odrębne składowe;
twierdzenie van Kampena daje
iloczyn wolny ich grup fundamentalnych
z amalgamacją nad $\pi_1(T_e)$.
W ten sposób wykonujemy $s-1$
amalgamacji i otrzymujemy przestrzeń
spójną, w której zachowane są obie
strony każdej pozostałej krawędzi.

Dla każdej krawędzi $e\notin T_\Gamma$
przyklejamy cylinder
$T_e\times[0,1]$, którego oba
końce leżą już w otrzymanej
przestrzeni spójnej. Jest to
przypadek HNN z
Twierdzenia~\ref{thm:tg-vk-amalgam-hnn}:
otrzymujemy nową stabilną literę
$t_e$ i relacje utożsamiające
obrazy elementów $G_e$ po obu
stronach. Pozostałych krawędzi jest
$m-(s-1)=m-s+1$, czyli dokładnie
$b_1(\Gamma)$.

Dla każdej krawędzi wybierzmy orientację oraz
oznaczmy monomorfizmy przy jej początku i końcu
przez $\iota_{e,-}$ i $\iota_{e,+}$.
Po zgodnym wyborze dróg bazowych końcowa
prezentacja jest ilorazem
\[
\left(*_{v\in V(\Gamma)}G_v\right)
 * F(\{t_e:e\notin T_\Gamma\}),
\]
gdzie $F$ oznacza grupę wolną, przez relacje
\[
\begin{cases}
\iota_{e,-}(a)=\iota_{e,+}(a),&e\in T_\Gamma,\\
t_e\iota_{e,-}(a)t_e^{-1}=\iota_{e,+}(a),
&e\notin T_\Gamma,
\end{cases}
\qquad a\in G_e.
\]
Są to relacje grupy fundamentalnej grafu grup
$\mathcal G$ z wybranym drzewem rozpinającym.
Konstrukcja topologiczna daje zatem wskazany
izomorfizm i pokazuje, że zmiana drzewa nie
zmienia typu izomorfizmu otrzymanej grupy.
\end{proof}

\begin{stwierdzenie}[Homologia grafu rozkładu]
\label{st:tg-jsj-h1}
Przy założeniach twierdzenia
\ref{thm:tg-jsj-graph-of-groups} oznaczmy
$B_v=H_1(M_v;\mathbb Z)$ i $B_e=H_1(T_e;\mathbb Z)$.
Wybieramy po jednej orientacji każdej krawędzi
geometrycznej grafu $\Gamma$ (także każdej pętli).
Wszystkie sumy po $e$ w tym stwierdzeniu
i jego dowodzie dotyczą właśnie tych krawędzi,
a nie obu ich orientacji.
Niech $j_{e,-}$ i $j_{e,+}$ będą homomorfizmami
$B_e\to\bigoplus_v B_v$ indukowanymi przez
oba włączenia brzegowe i włączenia odpowiednich
składników sumy prostej. Definiujemy
\[
D:\bigoplus_e B_e\longrightarrow\bigoplus_v B_v,
\qquad
D((a_e)_e)=\sum_e(j_{e,+}(a_e)-j_{e,-}(a_e)).
\]
Wtedy
\[
H_1(M;\mathbb Z)\cong
\operatorname{coker}D\oplus\mathbb Z^{b_1(\Gamma)}.
\]
Pierwszy składnik jest obrazem sumy homologii
części; rozszczepienie zależy od wyborów
w prezentacji grafu grup.
\end{stwierdzenie}
\begin{proof}
Abelianizacja iloczynu wolnego jest sumą prostą
abelianizacji czynników. W prezentacji z dowodu
twierdzenia \ref{thm:tg-jsj-graph-of-groups} każda
relacja krawędziowa po abelianizacji przyjmuje
postać $j_{e,+}(a)-j_{e,-}(a)=0$,
niezależnie od przynależności $e$ do drzewa:
wyrazy $t_e$ i $t_e^{-1}$ znoszą się.
Obraz wszystkich tych relacji w sumie prostej
jest dokładnie $\operatorname{im}D$,
gdyż homomorfizmy krawędziowe przechodzą przez
abelianizacje $G_e$. Żadna z relacji nie
zawiera niezerowej współrzędnej w grupie
swobodnej abelowej generowanej przez $t_e$.
Iloraz jest więc sumą prostą kokernu $D$
i grupy swobodnej abelowej rzędu $m-s+1$.
Stosujemy twierdzenie Hurewicza w wymiarze jeden.
Homomorfizmy $j_{e,\pm}$ na homologii nie
muszą być injektywne, choć odpowiednie
homomorfizmy na grupach podstawowych są
monomorfizmami.
\end{proof}

\begin{uwaga}[Drugi dowód przez ciąg Mayera--Vietorisa]
Wynik nie wymaga nieściśliwości torusów, jeżeli
rozważamy dowolny rozkład wzdłuż rozłącznych
dwustronnych torusów. W grafie przestrzeni
z dowodu twierdzenia \ref{thm:tg-jsj-graph-of-groups}
bierzemy otwarte pokrycie $U,V$: zbiór $U$
obejmuje części wierzchołkowe i końcowe odcinki
cylindrów, a $V$ ich środkowe odcinki.
Wybieramy odcinki tak, by ich przecięcie miało
po dwie składowe dla każdego cylindra.
Zbiór $U$ deformacyjnie ściąga się na sumę
rozłączną części $M_v$, zbiór $V$ na sumę
rozłączną torusów $T_e$, a $U\cap V$ na sumę
rozłączną dwóch kopii każdego $T_e$.
W stopniu jeden mapa Mayera--Vietorisa ma postać
\[
((a_e,b_e))_e\longmapsto
\left(\sum_e(j_{e,-}(a_e)+j_{e,+}(b_e)),
(-a_e-b_e)_e\right).
\]
W kokernie eliminujemy składnik homologii $V$:
relacja dla $(a_e,0)$ utożsamia klasę z tego
składnika z jej obrazem przez $j_{e,-}$,
a relacja dla $(0,b_e)$ narzuca różnicę
$j_{e,+}(b_e)-j_{e,-}(b_e)=0$.
Pozostaje dokładnie $\operatorname{coker}D$.

W stopniu zero analogiczna mapa wysyła
parę liczb $(a_e,b_e)$ na sumę ich wkładów
przy odpowiednich wierzchołkach oraz na
$-a_e-b_e$ przy krawędzi $e$. W jądrze mamy
$b_e=-a_e$ dla każdego $e$, a warunek
zanikania współrzędnych wierzchołkowych
jest równoważny zerowaniu brzegu całkowitego
1-łańcucha grafu o współczynnikach $b_e$.
Jądro jest więc izomorficzne z $H_1(\Gamma;\mathbb Z)$,
także dla pętli, których brzeg jest zerowy.
Z dokładności otrzymujemy ciąg
\[
0\to\operatorname{coker}D\to H_1(M;\mathbb Z)
\to H_1(\Gamma;\mathbb Z)\to0.
\]
Prawa grupa jest wolna abelowa rzędu $b_1(\Gamma)$,
więc ciąg rozszczepia się: wybieramy podniesienia
elementów jej bazy i rozszerzamy liniowo.
Dla rodziny pustej wynik jest bezpośredni,
bo $M$ jest jedyną częścią wierzchołkową.
\end{uwaga}

\begin{przyklad}[Torus nierozdzielający w wiązce nad okręgiem]
\label{ex:tg-jsj-sol-torus-bundle}
Niech
$A=\begin{psmallmatrix}2&1\\1&1\end{psmallmatrix}
\in SL_2(\mathbb Z)$ i rozważmy
trójrozmaitość będącą torusem odwzorowania
$A\colon T^2\to T^2$:
\[
M_A=(T^2\times[0,1])/
\bigl((x,1)\sim(Ax,0)\bigr).
\]
Macierz $A$ ma wartości własne
$(3\pm\sqrt5)/2$, więc jest
hiperboliczna. Rozmaitość $M_A$
ma geometrię typu Sol i
nie jest rozmaitością Seiferta.
Włókno $T^2$ jest torusem nierozdzielającym:
rozcięcie wzdłuż niego daje spójny iloczyn
$T^2\times[0,1]$. Z długiego ciągu grup homotopii
wiązki $T^2\to M_A\to S^1$ i równości
$\pi_2(S^1)=0$ wynika injektywność włączenia
grupy fundamentalnej włókna.
Torus ten jest więc nieściśliwy.
Wybierając orientację generatora grupy
fundamentalnej okręgu bazowego tak,
aby sprzężenie działało na $\pi_1(T^2)$
przez $A$ (odwrotny wybór daje $A^{-1}$),
otrzymujemy rozszerzenie HNN
\[
\pi_1(M_A)\cong
\langle x,y,t\mid
[x,y]=1,\;
txt^{-1}=x^2y,\;
tyt^{-1}=xy\rangle
\cong\mathbb Z^2\rtimes_A\mathbb Z.
\]
W tej prezentacji wektory zapisujemy kolumnowo:
obrazy $x,y$ odpowiadają pierwszej i drugiej
kolumnie macierzy $A$. Homomorfizm do
$\mathbb Z$ zadany przez $t\mapsto1$ i
$x,y\mapsto0$ ma przekrój będący homomorfizmem
$n\mapsto t^n$. Mnożenie w iloczynie półprostym
ma postać $(u,n)(v,m)=(u+A^nv,n+m)$.
Każdy element ma jednoznaczną postać $x^a y^b t^n$,
gdyż powyższe relacje pozwalają go do niej sprowadzić,
a w iloczynie półprostym różne trójki $(a,b,n)$
dają różne elementy.

Korzystamy z klasyfikacji geometrii wiązek torusowych
z monodromią Anosowa (Aschenbrenner--Friedl--Wilton,
\emph{3-Manifold Groups}, część o rozmaitościach Sol).
Wiązka jest nieprzywiedlna, ma geometrię Sol
i nie jest rozmaitością Seiferta. Nie jest też
atoroidalna, bo zawiera nieściśliwe włókno i nie ma
brzegu. Po rozcięciu wzdłuż włókna dostajemy
rozmaitość Seiferta $T^2\times[0,1]$ (okręgi
jednego czynnika torusa są włóknami).
Rodzina złożona z jednego torusa spełnia więc
warunki twierdzenia \ref{thm:tg-jsj-theorem},
a rodzina pusta ich nie spełnia. Jednoznaczność
w tym twierdzeniu identyfikuje włókno z torusem JSJ.
Graf ma jeden wierzchołek i jedną pętlę.
Nie musi zatem być drzewem; torus nierozdzielający
może dać pętlę i literę stabilną. Rachunek grupy
nie wymaga klasyfikacji geometrii.
\end{przyklad}

\begin{stwierdzenie}[Pierwsza homologia wiązki torusowej]
\label{st:tg-jsj-mapping-torus-h1}
Niech $A\in SL_2(\mathbb Z)$ i niech
$M_A=(T^2\times[0,1])/\bigl((x,1)\sim(Ax,0)\bigr)$
będzie torusem odwzorowania automorfizmu
$A\colon T^2\to T^2$.
Wówczas
\[
H_1(M_A;\mathbb Z)
\cong\mathbb Z\oplus
\operatorname{coker}(A-I_2:\mathbb Z^2\to\mathbb Z^2).
\]
Dla macierzy
$A=\begin{psmallmatrix}2&1\\1&1\end{psmallmatrix}$
z poprzedniego przykładu mamy
$H_1(M_A;\mathbb Z)\cong\mathbb Z$.
\end{stwierdzenie}

\begin{proof}
Jak w przykładzie \ref{ex:tg-jsj-sol-torus-bundle},
stosujemy twierdzenie \ref{thm:tg-vk-amalgam-hnn}
do cylindra torusa. Obie mapy na grupach są
izomorfizmami (jedna jest identycznością, druga
jest indukowana przez $A$). Dla dowolnego
$A\in SL_2(\mathbb Z)$ otrzymujemy więc
$\pi_1(M_A)\cong\mathbb Z^2\rtimes_A\mathbb Z$.
Oznaczmy generatory $\mathbb Z^2$ przez
$x,y$, a generator drugiego czynnika przez $t$.
Po abelianizacji obie relacje
$tut^{-1}=A(u)$ dla $u\in\mathbb Z^2$
przechodzą w $(A-I_2)u=0$.
Nie pojawia się żadna relacja na klasie
generatora $t$, więc
\[
\pi_1(M_A)_{\mathrm{ab}}
\cong\mathbb Z\oplus
\bigl(\mathbb Z^2/(A-I_2)\mathbb Z^2\bigr).
\]
Twierdzenie Hurewicza w wymiarze jeden
identyfikuje tę grupę z $H_1(M_A;\mathbb Z)$.

Dla wskazanej macierzy
$A-I_2=\begin{psmallmatrix}1&1\\1&0\end{psmallmatrix}$
ma wyznacznik $-1$, jest więc
automorfizmem grupy $\mathbb Z^2$.
Jej kokernel jest zerowy, co daje wynik.
\end{proof}

\begin{uwaga}[Rząd i torsja]
Ponieważ $\det A=1$, rozwinięcie wyznacznika
daje $\det(A-I_2)=2-\operatorname{tr}A$.
Jeżeli $\operatorname{tr}A\ne2$, to
$c=|2-\operatorname{tr}A|>0$. Niech $g$ będzie
dodatnim największym wspólnym dzielnikiem
czterech wpisów $A-I_2$. Z postaci normalnej
Smitha (pierwszy dzielnik determinantalny to $g$,
drugi to $c$) wynika
\[
H_1(M_A;\mathbb Z)\cong
\mathbb Z\oplus\mathbb Z/g\mathbb Z
\oplus\mathbb Z/(c/g)\mathbb Z .
\]
Mamy $g^2\mid c$, gdyż wszystkie wpisy są
podzielne przez $g$. Dla śladu $2$ rozróżniamy
$A=I_2$, kiedy $H_1\cong\mathbb Z^3$, oraz
$A\ne I_2$, kiedy $A-I_2$ ma rangę jeden
i $H_1\cong\mathbb Z^2\oplus\mathbb Z/g\mathbb Z$.
Na przykład dla $A=-I_2$ otrzymujemy
$\mathbb Z\oplus(\mathbb Z/2\mathbb Z)^2$,
a dla $A=\begin{psmallmatrix}1&k\\0&1\end{psmallmatrix}$
z $k\ne0$ otrzymujemy
$\mathbb Z^2\oplus\mathbb Z/|k|\mathbb Z$.
Wzory dotyczą wszystkich monodromii, nie tylko
hiperbolicznych.
\end{uwaga}

\begin{przyklad}[Ten sam graf JSJ, różna homologia]
\label{ex:tg-jsj-same-graph-different-h1}
Rozważmy dwie macierze
\[
A=\begin{pmatrix}2&1\\1&1\end{pmatrix},
\qquad
B=\begin{pmatrix}3&2\\1&1\end{pmatrix}
\]
z $SL_2(\mathbb Z)$.
Mają one odpowiednio ślady $3$ i $4$,
a obie mają wyznacznik równy $1$.
Ich monodromie są hiperboliczne,
więc odpowiadające im wiązki torusowe
$M_A$ oraz $M_B$ mają geometrię Sol.
W obu przypadkach rozkład JSJ
daje graf z jednym wierzchołkiem
i jedną krawędzią pętlową.

Homologie pierwsze są jednak różne.
Dla pierwszej macierzy $\det(A-I_2)=-1$,
stąd $H_1(M_A;\mathbb Z)\cong\mathbb Z$.
Dla drugiej mamy
\[
B-I_2=\begin{pmatrix}2&2\\1&0\end{pmatrix},
\qquad
\det(B-I_2)=-2.
\]
Ponieważ największy wspólny dzielnik
współczynników $B-I_2$ wynosi $1$,
jej postać normalna Smitha ma
niezerowe współczynniki diagonalne
$1$ i $2$. Poprzednie stwierdzenie daje
\[
H_1(M_B;\mathbb Z)\cong
\mathbb Z\oplus\mathbb Z/2\mathbb Z.
\]
Przestrzenie $M_A$ i $M_B$ nie są
zatem homeomorficzne, mimo że
ich nieoznakowane grafy JSJ są identyczne.
Przykład dowodzi konkretnie,
że sam graf dualny rozkładu
nie wystarcza do odtworzenia
trójrozmaitości: trzeba zachować
dane monodromii i sklejania.
\end{przyklad}

\section{Przykłady: węzły hiperboliczne, torusowe i satelitarne}
\label{sec:tg-jsj-knots}

\begin{stwierdzenie}[Graf JSJ zewnętrza węzła jest drzewem]
\label{st:tg-jsj-knot-tree}
Dla każdego węzła $K\subset S^3$ graf
$\Gamma_{\mathrm{JSJ}}(X_K)$ jest drzewem.
\end{stwierdzenie}
\begin{proof}
Zewnętrze $X_K$ jest nieprzywiedlne: każda gładko
zanurzona sfera w jego wnętrzu rozdziela $S^3$
na dwie kule na mocy twierdzenia Schoenfliesa
w kategorii gładkiej. Spójne otoczenie $N(K)$
jest rozłączne ze sferą i leży po jednej stronie,
więc druga kula zawiera się w $X_K$.
Zatem stosuje się twierdzenie \ref{thm:tg-jsj-theorem}.

Każdy torus $T_e$ rodziny JSJ rozdziela $S^3$:
z dualności Aleksandra mamy
$\widetilde H_0(S^3\setminus T_e;\mathbb Z)
\cong H^2(T_e;\mathbb Z)\cong\mathbb Z$,
więc dopełnienie ma dokładnie dwie składowe.
Obie zawierają punkty dowolnie blisko torusa
po odpowiednich stronach jego kołnierza.
Ponieważ torus leży we wnętrzu $X_K$, punkty te
można wybrać w $X_K\setminus T_e$.
Nie mogą być połączone ścieżką w tym zbiorze,
gdyż nie są połączone w $S^3\setminus T_e$.
Torus rozdziela więc także $X_K$.
Ze stwierdzenia \ref{st:tg-jsj-bridge} każda
krawędź grafu jest mostem, a graf jest spójny.
Spójny multigraf, którego każda krawędź jest
mostem, nie ma cykli (w tym pętli ani cykli
złożonych z dwóch krawędzi równoległych).
Jest zatem drzewem.
\end{proof}

\begin{uwaga}[Dlaczego homologia grafu JSJ znika]
Dla zewnętrza węzła poprzedni wynik można również
uzyskać bezpośrednio z homologii. Istotnie, z rozkładu
w stwierdzeniu \ref{st:tg-jsj-h1} grupa
$\mathbb Z^{b_1(\Gamma)}$ jest składnikiem prostym
$H_1(X_K;\mathbb Z)\cong\mathbb Z$.
Samo to ograniczenie daje jedynie $b_1(\Gamma)\leq1$.
Ponieważ meridian brzegu $X_K$ leży w jednej z części
rozkładu, jego klasa należy do pierwszego składnika,
czyli obrazu sumy homologii części. Jest ona generatorem
całej grupy na mocy stwierdzenia
\ref{thm:kt-peripheral-homology}. Pierwszy składnik
jest więc całą grupą, a $b_1(\Gamma)=0$.
Ze spójności grafu i wzoru $b_1=m-n+1$
z twierdzenia \ref{thm:tg-wedge} ponownie wynika,
że graf JSJ jest drzewem.
\end{uwaga}

\begin{przyklad}[Trzy różne geometrie, różne dane grafowe]
Niech $K\subset S^3$ będzie węzłem, a
$X_K=S^3\setminus\operatorname{int}N(K)$ jego
zewnętrzem. Nieprzywiedlność $X_K$ uzasadniliśmy
w dowodzie stwierdzenia \ref{st:tg-jsj-knot-tree}.
Dla węzła ósemkowego wnętrze $X_K$ jest
hiperboliczne o skończonej objętości, a samo
$X_K$ jest atoroidalne (por. sekcję
\ref{sec:kt-hyperbolic});
jego graf JSJ ma jeden wierzchołek i żadnej krawędzi.
Dla nietrywialnego węzła torusowego $T(p,q)$
zewnętrze jest rozmaitością Seiferta i również
ma graf JSJ z jednym wierzchołkiem. Sam graf nie
rozróżnia więc obu przypadków: potrzebna jest
etykieta rodzaju części oraz jej topologia.

Węzeł satelitarny ma w swoim zewnętrzu torus nieściśliwy,
który nie jest równoległy do toroidalnego brzegu.
Rozważmy konkretnie kabel węzła hiperbolicznego
$J$: w torusie $N(J)$ wybieramy prostą krzywą
klasy $p\mu_J+q\lambda_J$ z $\gcd(p,q)=1$ i
$|q|\ge2$, po czym odsuwamy ją do wnętrza.
Niech $K$ będzie tym kablem, a
$P=N(J)\setminus\operatorname{int}N(K)$.
Korzystamy ze standardowego opisu przestrzeni
kablowej jako rozmaitości Seiferta nad pierścieniem
z jednym włóknem wyjątkowym rzędu $|q|$
i z twierdzenia o torusie towarzysza w rozkładzie
JSJ kabla (R. Budney, \emph{JSJ-decompositions of
knot and link complements in the 3-sphere},
arXiv:math/0506523).
Torus $\partial N(J)$ jest jedynym torusem JSJ:
oddziela część kablową $P$ od hiperbolicznego
zewnętrza $X_J$, które nie ma własnych torusów JSJ.
Graf ma więc dwa wierzchołki i jedną krawędź.
Stosując twierdzenie \ref{thm:tg-jsj-graph-of-groups}
do tego drzewa, dostajemy przy odpowiednich
punktach bazowych
$\pi_1(X_K)\cong\pi_1(P)*_{\mathbb Z^2}\pi_1(X_J)$.
Torus krawędziowy wnosi dwa generatory, meridian i długość,
których obrazy po obu stronach sklejania zależą od typu kabla. Wyróżnienie grup obwodowych
$\pi_1(T^2)\cong\mathbb Z^2$ jest konieczne,
aby odzyskać sposób sklejenia tych części.
\end{przyklad}

\begin{uwaga}[Graf a geometria trójrozmaitości]
Rozkład JSJ ma szczególne znaczenie po
geometryzacji Thurstona. Części atoroidalne,
które nie są rozmaitościami Seiferta, mają
w rozkładzie JSJ zwartej, orientowalnej,
nieprzywiedlnej rozmaitości o pustym lub
torusowym brzegu, na mocy geometryzacji
wnętrza z zupełną metryką hiperboliczną
o skończonej objętości; korzystamy tu z
twierdzenia 1.7.6 z przywołanej monografii
\emph{3-Manifold Groups}.
Graf dualny porządkuje części, zaś grupy,
parametry torusów brzegowych i homeomorfizmy
sklejania zapisują dodatkowe dane topologiczne.
To właśnie te dane łączą teorię grafów z
chirurgią Dehna i rachunkiem Kirby'ego. Sama
izomorficzna struktura abstrakcyjnych grup na wierzchołkach
i krawędziach nie odzyskuje automatycznie pełnych map
sklejania torusów; do rekonstrukcji geometrycznej
należy zachować dane peryferyjne i ich utożsamienia.
\end{uwaga}

\begin{uwaga}[Literatura]
Klasyczne twierdzenie o rozkładzie pochodzi z prac
W. Jaco, P. Shalena oraz K. Johannsona.
Dalsze przykłady i strukturę grup fundamentalnych
można znaleźć w monografii M. Aschenbrennera,
S. Friedl, H. Wilton, \emph{3-Manifold Groups}.
\end{uwaga}


